\documentclass[12pt,a4paper]{article}
\usepackage[utf8]{inputenc}
\usepackage[vietnam]{babel}
\usepackage[T1]{fontenc}
\usepackage{geometry}
\geometry{a4paper, margin=1in}
\usepackage{hyperref}
\usepackage{listings}
\usepackage{xcolor}

% Cấu hình màu sắc cho code
\definecolor{codegreen}{rgb}{0,0.6,0}
\definecolor{codegray}{rgb}{0.5,0.5,0.5}
\definecolor{codepurple}{rgb}{0.58,0,0.82}
\definecolor{backcolour}{rgb}{0.95,0.95,0.92}

\lstdefinestyle{mystyle}{
    backgroundcolor=\color{backcolour},   
    commentstyle=\color{codegreen},
    keywordstyle=\color{blue},
    numberstyle=\tiny\color{codegray},
    stringstyle=\color{codepurple},
    basicstyle=\ttfamily\footnotesize,
    breakatwhitespace=false,         
    breaklines=true,                 
    captionpos=b,                    
    keepspaces=true,                 
    numbers=left,                    
    numbersep=5pt,                  
    showspaces=false,                
    showstringspaces=false,
    showtabs=false,                  
    tabsize=2
}
\lstset{style=mystyle}

\title{\textbf{Tài liệu bàn giao và Hướng dẫn vận hành:\\ Multi-WAN Packet Forwarder}}
\author{Hoàng Bửu}
\date{\today}

\begin{document}

\maketitle

\section{Tổng quan dự án}
Dự án \textbf{Multi-WAN Packet Forwarder} là giải pháp dataplane chạy trên tầng người dùng (user-space), sử dụng cơ chế \texttt{AF\_PACKET (TPACKET\_V3)}. Mục tiêu chính là xử lý lưu lượng IP tốc độ cao: tiếp nhận, sửa đổi header, phân mảnh gói tin và cân bằng tải qua các đường truyền WAN/Tunnel.

\subsection{Đặc tính kỹ thuật}
\begin{itemize}
    \item \textbf{Tối ưu hóa hiệu năng:} Xử lý gói tin trực tiếp qua memory-mapped ring, giảm thiểu overhead do chuyển đổi ngữ cảnh giữa Kernel và User-space.
    \item \textbf{Xử lý gói tin:} Hỗ trợ phân mảnh (Fragmentation) động dựa trên MTU của Tunnel. Sử dụng thuật toán RFC 1624 để tính toán lại Checksum nhanh.
    \item \textbf{Cân bằng tải:} Phân phối lưu lượng dựa trên Hash 5-tuple, đảm bảo tính toàn vẹn thứ tự gói tin cho mỗi luồng (flow).
    \item \textbf{Quản lý tập trung:} Cấu hình thiết bị (MAC, Interface, Routing) được quản lý qua cơ sở dữ liệu PostgreSQL.
\end{itemize}

\section{Kiến trúc hệ thống}
Người kế nhiệm cần nắm rõ hai thành phần vận hành song song:
\begin{enumerate}
    \item \textbf{Control Plane (\texttt{main.c}):} Đóng vai trò thực thi lệnh điều khiển. Nó lắng nghe kết nối TCP, truy vấn cấu hình từ DB và điều phối tài nguyên cho Dataplane.
    \item \textbf{Data Plane (\texttt{afpkt.c}):} Tập hợp các luồng Omni-Worker chạy theo mô hình \textit{run-to-completion}. Mỗi worker được gán cứng vào một lõi CPU (Core Pinning) để tránh tranh chấp tài nguyên.
\end{enumerate}

\section{Yêu cầu môi trường}
Trước khi triển khai, hệ thống cần đáp ứng:
\begin{itemize}
    \item \textbf{Kernel:} Hỗ trợ \texttt{TPACKET\_V3} và \texttt{PACKET\_FANOUT}.
    \item \textbf{Thư viện:} Cài đặt \texttt{libpq-dev} để kết nối PostgreSQL.
    \item \textbf{Quyền hạn:} Mọi thao tác vận hành yêu cầu quyền \texttt{root}.
\end{itemize}

\section{Hướng dẫn triển khai (Deployment)}

\subsection{Biên dịch mã nguồn}
Sử dụng \texttt{make} để đóng gói ứng dụng:
\begin{lstlisting}[language=bash]
# Cài đặt thư viện phát triển PostgreSQL
sudo apt-get update && sudo apt-get install libpq-dev -y

# Biên dịch dự án
make clean && make
\end{lstlisting}
Sau khi biên dịch thành công, file thực thi \texttt{mwanc} sẽ xuất hiện trong thư mục gốc.

\subsection{Khởi tạo cơ sở dữ liệu}
Dữ liệu cấu hình mạng phải được nạp trước vào PostgreSQL:
\begin{lstlisting}[language=bash]
psql -U [user] -d [database] -f test_connect_db/init_db.sql
\end{lstlisting}

\section{Vận hành ứng dụng}

\subsection{Bước 1: Khởi chạy Daemon}
Kích hoạt ứng dụng và kết nối tới DB:
\begin{lstlisting}[language=bash]
sudo ./mwanc --db-host 127.0.0.1 \
             --db-port 5432 \
             --db-user postgres \
             --db-name mydb \
             --listen-port 8080
\end{lstlisting}

\subsection{Bước 2: Nạp cấu hình từ Backend}
Để kích hoạt dataplane cho một node cụ thể (ví dụ \texttt{server1}), sử dụng script giả lập:
\begin{lstlisting}[language=bash]
# Cú pháp: ./script [node_id] [port]
./test_connect_db/mock_backend.sh server1 8080
\end{lstlisting}

\section{Cấu trúc mã nguồn (Dành cho nhà phát triển)}
Khi cần bảo trì hoặc nâng cấp, hãy lưu ý các thư mục sau:
\begin{itemize}
    \item \texttt{src/proto/}: Chứa logic quan trọng nhất về phân mảnh gói tin và tính toán checksum.
    \item \texttt{src/userio/}: Quản lý vòng lặp thu thập gói tin (Packet RX/TX) hiệu suất cao.
    \item \texttt{src/system/}: Các hàm can thiệp hệ thống như tắt tính năng Offload của card mạng (TSO, GSO) để tránh xung đột với User-space Dataplane.
    \item \texttt{src/config/}: Logic giao tiếp và parse dữ liệu từ PostgreSQL.
\end{itemize}

---
\textit{Tài liệu này được biên soạn để phục vụ việc bàn giao. Mọi thay đổi về kiến trúc sau này cần được cập nhật vào mục 2 và 6.}

\end{document}